\documentclass{exam}
\usepackage{graphicx} % Required for inserting images
\usepackage{algorithmicx}
\usepackage{algpseudocode}
\usepackage{geometry}[border=1in]
\usepackage{algorithm}
\usepackage{amsmath}
\usepackage{amssymb}
\usepackage{amsthm}
\usepackage{listings}
\usepackage{mathtools}
\usepackage{hyperref}
\usepackage{framed}
\DeclarePairedDelimiter\ceil{\lceil}{\rceil}
\DeclarePairedDelimiter\floor{\lfloor}{\rfloor}

\newtheoremstyle{prob}% name 
{7pt}% Space above 
{7pt}% 
{}% 
{0pt}%Indent amount 1 
{\bfseries}% Theorem head font 
{.}% 
{\newline}% 
{}%

\theoremstyle{prob}
\newtheorem*{probstate}{PROBLEM}

\printanswers

\title{Homework 4: Greedy Algorithms and SSSP}
\author{COMP221 Spring 2026 - Suhas Arehalli}
\date{}

\begin{document}

\maketitle

Complete the problems below. Check the course website \& syllabus for further instructions.

If any problem is unclear, or you think you found a typo, please let me know ASAP so I can clarify!

\section*{Problems}

\begin{questions}
    \question \textbf{Other Kinds of Greed} (\textit{Based on Erickson Ch.4 Prob. 1})

    Recall that the greedy scheduling problem in Erickson 4.2 is the defined (somewhat informally) as follows: 

    \begin{probstate}[Scheduling]
        \textbf{Input:} A set of $n$ events with start/end times $E = \{(s_1, e_1), \dots, (s_n, e_n)\}$ \\
        \textbf{Output:} The largest $X \subseteq E$ such that no events in $X$ overlap.
    \end{probstate}

    For each of the following greedy strategies to build a schedule, either prove that the strategy works via exchange argument or prove the strategy does not work via counterexample. You may assume ties can be broken either way.

    \begin{parts}
        \part[6] (Repeatedly) select the event with the latest end time that doesn't conflict with events we've chosen so far.

        \part[6] (Repeatedly) select the event with the latest start time that doesn't conflict with events we've chosen so far.
        
        \part[6] (Repeatedly) select the event with the shortest duration that doesn't conflict with events we've chosen so far. 

        \part[2] (Repeatedly) select the event that doesn't conflict with events we've chosen so far and has the fewest conflicts with yet-to-be-chosen events. 

        \textbf{**HINT**:} Despite appearances, this strategy does \textit{not} find optimal schedules. Find a counterexample!

    \end{parts}

    \question \textbf{Maintaining Relaxation}

    You can compute shortest path distances sing Dijkstra's (or Bellman-Ford, etc.), but these algorithms take a not-insignificant amount of time (an optimized version of Dijkstra's takes $\Theta(E + V\log V)$ time with no negative edge weights, and less optimized versions like ours are are worse). This means that in practice, we don't want to have to recompute shortest paths when we don't actually need to. This is particularly true when we may consider re-using our shortest path algorithm over a graph that is only minimally different from one which we've already computed shortest paths over! 

    Suppose you are given a (directed) Graph $G = (V, E)$, a weight function $W: E \to R$ (i.e., non-negative weights), a source vertex $s \in V$. You may assume there are no negative-weight cycles.

    \begin{parts}
        \part[8] Suppose you are also given a map $d$ where $d(v) = \delta(s,v)$. Then an edge $e \in E$ decreases weight: You are given $w': E \to \mathbb{R}$ such that for some $c \geq 0$ \textbf{and} the edge $e$ that has been modified.
        \begin{align*}
            w'(u,v) = \begin{cases}
                w(u,v), & (u,v) \neq e \\
                w(u,v) - c, &(u,v) = e
            \end{cases}
        \end{align*}
        Construct an algorithm $\Call{StillShortest}{G, d, w', e}$ that runs in $\Theta(1)$ time to determine if $d$ still contains shortest paths for $G$ under $w'$ (rather than $w$). \label{4:decrease-weight}


        \part[2] Write a (short) proof of correctness for your algorithm. You should reference results we've seen in class/the textbook, but make sure you understand when they apply!


        \part[3] Consider a different scenario: Suppose add a new edge $(u,v)$ to $E$ for $u,v \in V$ and assign it a new weight $w(e) = c$. Write an algorithm to determine whether we need to update our shortest path distances $d$.

       \textbf{ **HINT**:} For any graph $G = (V, E)$ with weights $w: E \to \mathbb{R}$, consider the complete graph $G_c = (V, V^2)$ where 

        \begin{align*}
            w_c(u,v) = \begin{cases}
                w(u,v), &(u,v) \in E \\
                \infty, &(u,v) \notin E
            \end{cases}
        \end{align*}

        Are shortest paths any different between $G_c$ with $w_c$ and $G$ with $w$? With that in mind, how does the scenario in this part compare to (\ref{4:decrease-weight})?


        \part[2] Suppose we have the weights of some edges modified (either increased or decreased!), but this time we maintain not just $d: V \to \mathbb{R} \cup \{\infty\}$, but $\pi: V \setminus \{s\} \to V$, which maps a vertex to it's parent in the shortest path tree in the original graph (pre-modification!). Note (1) what can go wrong if we \textit{increase} the weights of graphs and (2) describe how $\pi$ solves this problem. 


        \part[3] Now consider a new (but related!) problem: Suppose we are given a graph $G=(V,E)$, $w: E \to \mathbb{R}$, and a parent function $\pi: V \setminus \{s\} \to V$. Write an algorithm to determine whether $\pi$ defines a single-source shortest path tree rooted at $s$. 


        \part[2] What is the runtime of your algorithm? Remember to consider both the number of vertices and edges in the graph!


    \end{parts}

\end{questions}

\end{document}